\chapter{Evidence classes and the claim ledger}
\label{book:appendixC}

Every result in this book belongs to one of the evidence classes below. The classes are
not a ranking of worth. A clean proof and a careful measurement answer different
questions, and the most common error is to let evidence of one class stand in for
another.

\section{Eight classes of evidence}
\begin{center}\small
\begin{tabular}{@{}P{0.18\linewidth}P{0.27\linewidth}P{0.27\linewidth}P{0.20\linewidth}@{}}
\toprule
Class & What it can establish & What it cannot & Example in the book\\\midrule
Proof under premises & A property for every input satisfying stated assumptions &
That the assumptions hold for a real system & Scan associativity (Ch.~7, 26)\\
Finite reference test & That an implementation agrees with an oracle on declared cases &
Behaviour on untested inputs or other implementations & KV-cache parity (Ch.~35)\\
Negative control & That a test can detect a specific failure & That no other failure
exists & Future-copy predictor (Ch.~2)\\
Synthetic trace & Exact behaviour of a mechanism on constructed input & Usefulness on
natural data & Selective-memory traces (Ch.~19)\\
Numerical emulation & The effect of a numerical policy on a fixed model & Hardware
speed or storage savings & KV precision rollout (Ch.~37)\\
Trained held-out result & Predictive performance under a split and budget & Biological
mechanism; transfer beyond the split & Motif pilots (Ch.~22--23)\\
Hardware measurement & Time, memory or energy on stated hardware and workload &
Performance on other workloads or devices & None reported in this edition\\
Biological observation & What happened in a declared experiment & Generality beyond its
conditions & Cited primary studies only\\
\bottomrule
\end{tabular}
\end{center}
A ninth category is the \emph{hypothesis}: a falsifiable statement with no supporting
evidence yet. Hypotheses are valuable, and the book states many of them, but they are
never reported as results.

\section{A claim template}
Chapters~29, 51 and~52 use a ledger that binds a claim to the evidence that would
support or refute it. A minimal entry has these fields:
\begin{center}\small
\begin{tabular}{@{}P{0.24\linewidth}P{0.70\linewidth}@{}}
\toprule
Field & Content\\\midrule
Claim & Exact wording, with no adjectives that the evidence does not carry\\
Scope & Task population, information access, budget and units\\
Evidence class & One of the classes above, and the artifacts that carry it\\
Baseline & The closest competitor, given the same opportunity\\
Falsifier & A result, fixed before the run, that would weaken or reject the claim\\
Uncertainty & Sampling unit, number of independent units, interval method\\
Artifacts & Content identities of code, data, model, environment and raw results\\
Status & Hypothesized, implemented, tested, benchmark-demonstrated or replicated\\
Failures & Linked counterexamples, kept after any resolution\\
\bottomrule
\end{tabular}
\end{center}

\section{Ten questions to ask of any result}
\begin{enumerate}
\setlength{\itemsep}{1pt}
\item Is the target aligned with the information the model may legally use?
\item Could overlap, homology, reverse complements or shared sources leak across the split?
\item Does a trivial baseline (composition, length, majority class) explain the score?
\item Were all runs, seeds and failed configurations reported, or only the best?
\item Is the uncertainty computed over the actual independent units?
\item Do the compared systems have matched data, compute, tuning and state budgets?
\item Does an ablation or sham intervention isolate the claimed mechanism?
\item Is a measured number labelled with its hardware, workload and timing boundary?
\item Could the result be reproduced from the recorded artifacts by someone else?
\item Is the conclusion worded no more strongly than its evidence class allows?
\end{enumerate}
A result that survives these questions is not thereby important. It has become
something that can be built on.
